\documentclass[aps,pra,twocolumn,superscriptaddress,longbibliography]{revtex4-2}

\usepackage{amsmath,amssymb,amsthm,graphicx,xcolor,bm,braket,dcolumn,booktabs}
\usepackage{hyperref}

\newtheorem{theorem}{Theorem}
\newtheorem{observation}{Observation}

\begin{document}

\title{The Three-Center Electron-Repulsion Integral Is an Elliptic
Bessel Moment:\\
Genus One at the Third Center in Momentum-Space Slater Theory}
\thanks{Paper 59 in the Geometric Vacuum series}

\author{J.~Loutey}
\affiliation{Independent Researcher, Kent, Washington}

\date{August 16, 2026}

%% ========================================================================
%% ABSTRACT
%% ========================================================================

\begin{abstract}
Two-center electron-repulsion integrals over Slater-type orbitals close
in elementary functions plus the exponential integral, the logarithm, and
Euler's~$\gamma$---a genus-zero transcendence class---which is why the
two-center problem has been solved since Roothaan and Ruedenberg.  The
genuine three-center integral has resisted a closed form for seventy years,
and the standard verdict is that it is simply hard: reduce it to an infinite
series or a one-dimensional quadrature and be done.  We give it a structural
reading instead.  Evaluating the two-electron three-center integral
$(XY|XZ)$ over $1s$ orbitals in momentum space---the Fock projection under
which Coulomb--Sturmians become hyperspherical harmonics, the setting native
to the Averys' generalized-Sturmian program---the three centers enter as
three \emph{phases} rather than three foci, so the coordinate obstruction
(prolate spheroidal coordinates admit exactly two foci) dissolves.
\textbf{[MEASURED]} The resulting momentum integral reproduces the exact
integral (agreement $1.9\times10^{-14}$ against a closed-form Gaussian
reference; $4.3\times10^{-7}$ against an independent Slater evaluation), and
its angular part collapses to a single spherical Bessel kernel.
\textbf{[MEASURED / classical]} What remains---the radial content---is a
two-scale Bessel moment whose period is a \emph{complete elliptic integral}
(verified to 31 digits), living on the genus-one curve
$y^{2}=(c_{1}k^{2}+1)(c_{2}k^{2}+1)$; it degenerates to the genus-zero
two-center case exactly on the diagonal $c_{1}=c_{2}$.  \textbf{[SYMBOLIC +
MEASURED]} The one-mass slice is the Laplace transform of the elliptic
differential---an elliptic deformation of the Bessel $K_{0}$, reducing to
$K_{0}(D)$ when the two orbital scales coincide and to the complete elliptic
integral $K$ at zero separation---and satisfies an explicit fourth-order
Picard--Fuchs equation with polynomial coefficients (residual
$\sim\!10^{-24}$), whose characteristic exponents are the four branch points
of the curve.  \textbf{[OBSERVATION]} This places a molecular integral inside
the modern theory of Bessel moments and two-mass elliptic Feynman integrals;
a targeted literature search finds that connection unmade in either field.
\textbf{[OPEN]} We do \emph{not} give the finite closed form, and we show the
natural route to it is obstructed: the elliptic-dilogarithm (unequal-mass
sunrise) mechanism does not engage, because the modulus Picard--Fuchs source
is not an elementary subtopology but lies in the integral's own Gauss--Manin
$D$-module---the branch-point boundary that would supply the sunrise's tadpole
vanishes identically.  \textbf{[OBSERVATION]} The object is therefore a
genuinely new elliptic transcendent one level \emph{above} the
elliptic-dilogarithm tower, a period of a rank-four irregular connection;
\textbf{[OPEN]} whether it admits any closed form is a differential-Galois
question we leave open.  The message is not a solution
but a relocation: the third center does not make the radial integral messier
within the two-center class, it raises the transcendence one genus, past the
reach of the sunrise machinery.
\end{abstract}

\maketitle

%% ========================================================================
\section{The two-center class, and the wall beyond it}
%% ========================================================================

The electron-repulsion integral over Slater-type orbitals on two centers is
one of the oldest solved problems in molecular quantum mechanics.  In prolate
spheroidal coordinates, with the two nuclei at the foci, the Neumann
expansion of $1/r_{12}$ terminates against the orbital angular momenta, and
every integral reduces to elementary functions together with the exponential
integral $E_{1}$, the logarithm, and Euler's constant~$\gamma$.  In the
transcendence bookkeeping of the Geometric Vacuum program (Paper~18), these
are \emph{embedding-tier} constants---the cost of the continuum two-center
geometry rather than content of the discrete graph---and, decisively, they
sit at transcendence \emph{weight one}: no $\pi$ survives (it cancels against
the Coulomb $4\pi$), and no dilogarithm appears, even though the ordered
radial integral has the simplex shape that generically produces one.  This
weight-one closure is what lets the two-center repulsion density be
\emph{decided} rather than merely counted (Paper~58): the closed form is a
rational combination of $\{1,\,e^{-\lambda R}\}$, and Lindemann--Weierstrass
separates its zeros.

The genuine three-center integral has no such story.  Prolate spheroidal
coordinates are built from exactly two foci; a third nucleus has nowhere to
sit.  Gaussians escape this because two Gaussians on different centers
multiply into a single Gaussian on a third point---Slater orbitals have no
product theorem---and this, not any Geometric-Vacuum-specific limitation, is
the seventy-year-old reason Gaussian bases displaced Slater orbitals for
polyatomic quantum chemistry.  The modern state of the art for three-center
Slater repulsion integrals (Ref.~\cite{ozdogan2012}) reduces them to a single
infinite expansion or a one-dimensional numerical quadrature; the one
celebrated closed form in correlated two-electron theory, Fromm and Hill's
analytic evaluation~\cite{fromm_hill1987}, is a \emph{dilogarithm}---a
genus-zero object---and belongs to a different integral topology (the
one-center, $r_{12}$-correlated master integral), which lacks the structure
we are about to exhibit.

This paper asks what the three-center integral \emph{is}, transcendentally,
and answers: it is an elliptic object.  The two-center weight-one closure is
a genus-zero fact and it stops at two centers.  The third center raises the
transcendence one genus, to the elliptic (genus-one) world of Bessel moments
and two-mass sunrise Feynman integrals.  We establish this constructively,
validate every step numerically, place the object in the literature, and mark
the finite closed form as open.

\emph{Lineage and scope.}  The momentum-space route is not new machinery.
Momentum-space Coulomb--Sturmians are hyperspherical harmonics under Fock's
conformal projection~\cite{fock1935}, the framework the Averys built into the
generalized-Sturmian method~\cite{avery1989,avery2006}, and the closest prior
art to the present calculation is the Averys' own momentum-space evaluation of
Coulomb--Sturmian repulsion integrals~\cite{avery_avery_4center}, which
stops---as all prior work does---at expansion and numerics.  What is new here
is the reading: that the residual radial object is a specific elliptic Bessel
moment, and hence that a molecular integral is an instance of the
elliptic-Feynman-integral class.  We work throughout with $1s$ orbitals at a
common exponent; the generalization to higher $(n,\ell,m)$ multiplies the
integrand by solid harmonics and shifts exponents but does not change the
genus, and we return to this in Sec.~\ref{sec:scope}.

%% ========================================================================
\section{The momentum-space reduction}
\label{sec:momentum}
%% ========================================================================

Write the chemist integral for two densities
$\rho_{1}=\chi_{X}\chi_{Y}$ and $\rho_{2}=\chi_{X}\chi_{Z}$, two two-center
$1s$ densities that share the center $X$ on different axes.  With the Fourier
representation of the kernel,
$1/|\bm r_{1}-\bm r_{2}| = (2\pi^{2})^{-1}\!\int d^{3}k\,
e^{i\bm k\cdot(\bm r_{1}-\bm r_{2})}/k^{2}$, the integral becomes
%
\begin{equation}
  (XY|XZ) = \frac{1}{2\pi^{2}}\int \frac{d^{3}k}{k^{2}}\,
  \tilde\rho_{1}(\bm k)\,\overline{\tilde\rho_{2}(\bm k)},
  \qquad
  \tilde\rho(\bm k)=\!\int\! \rho(\bm r)\,e^{i\bm k\cdot\bm r}\,d^{3}r .
  \label{eq:momentum}
\end{equation}
%
Each center now appears only through a phase $e^{i\bm k\cdot\bm R}$: three
centers are three phases, not three foci, and there is nothing to
diagonalize.  Equation~\eqref{eq:momentum} is exact; the only content is the
density Fourier transform $\tilde\rho$ and the $\bm k$ integral.

We verify Eq.~\eqref{eq:momentum} two independent ways on the reference
geometry $X=(0,0,0)$, $Y=(0,0,2)$, $Z=(1.5,0,0.5)$ (bohr), $\zeta=1$, for
which a McMurchie--Davidson evaluation over $8$-Gaussian $1s$ fits gives the
ground truth $(XY|XZ)=0.204941722$.  \textbf{[MEASURED]} Using the closed-form
Fourier transform of the Gaussian-fitted density (a product of Gaussians is a
Gaussian) in Eq.~\eqref{eq:momentum} reproduces the McMurchie--Davidson value
on the \emph{same} fits to $1.9\times10^{-14}$: the momentum formula, the
Fourier and Coulomb conventions, and the $\bm k$-quadrature are correct.
Independently, computing the \emph{true} Slater density transform (no Gaussian
fit) by a one-parameter reduction of the two-center Yukawa product reproduces
the ground truth to $4.3\times10^{-7}$.  Both are pinned in the regression
suite~\cite{test_routeC}.

\emph{The angular integral closes to a spherical Bessel kernel.}  Writing
each Slater factor through $e^{-\zeta r}=-\partial_{\zeta}(e^{-\zeta r}/r)$
and the Feynman parametrization of the resulting Yukawa product, the three
phases combine into a \emph{single} phase $e^{i\bm k\cdot\bm W}$ with
$\bm W(s,t)=(t-s)\bm X+s\bm Y-t\bm Z$ ($s,t$ the two Feynman parameters), so
the angular integral is elementary,
%
\begin{equation}
  \int d\Omega_{\bm k}\, e^{i\bm k\cdot\bm W}
  = 4\pi\, j_{0}(k|\bm W|)
  = 4\pi\,\frac{\sin(k|\bm W|)}{k|\bm W|}.
  \label{eq:angular}
\end{equation}
%
The whole integral reduces to a two-dimensional Feynman integral over a
one-dimensional radial ($k$) integral,
%
\begin{align}
  (XY|XZ) = \frac{8}{\pi}\,
  &\partial_{\zeta_{a}}\partial_{\zeta_{b}}
   \partial_{\zeta_{c}}\partial_{\zeta_{d}}
   \!\int_{0}^{1}\!\!ds\!\int_{0}^{1}\!\!dt\!\int_{0}^{\infty}\!\!dk\;
   j_{0}(k|\bm W|)\notag\\[-2pt]
  &\times\frac{e^{-D_{1}\Delta_{1}}}{\Delta_{1}}\,
          \frac{e^{-D_{2}\Delta_{2}}}{\Delta_{2}}
  \;\Bigg|_{\zeta=1},
  \label{eq:reduced}
\end{align}
%
with $\Delta_{1}=\sqrt{c_{1}k^{2}+1}$, $\Delta_{2}=\sqrt{c_{2}k^{2}+1}$,
$c_{1}=s(1-s)$, $c_{2}=t(1-t)$, and $D_{1}=|X-Y|$, $D_{2}=|X-Z|$.  The
$\partial_{\zeta}$'s (a bookkeeping device restoring the Slater from the
Yukawa) and the $(s,t)$ integrations are elementary operations; the
transcendental content is entirely in the radial factor.
Equation~\eqref{eq:reduced} reproduces the ground truth to
$1.8\times10^{-6}$ \textbf{[MEASURED]}, the residual being the finite-order
angular and Feynman quadrature.

%% ========================================================================
\section{The transcendence jump: genus zero to genus one}
\label{sec:genus}
%% ========================================================================

The radial factor of Eq.~\eqref{eq:reduced} carries two dispersion factors
$\Delta_{1},\Delta_{2}$ with \emph{different} scales $c_{1}\neq c_{2}$.  A
single such factor is benign.  Under Fock's substitution
$k\sqrt{c}=m\sinh\theta$ it linearizes, and \textbf{[MEASURED, $6\times
10^{-18}$]}
%
\begin{equation}
  \int_{0}^{\infty}\!\!\cos(kb)\,
  \frac{e^{-D\sqrt{c\,k^{2}+m^{2}}}}{\sqrt{c\,k^{2}+m^{2}}}\,dk
  = \frac{1}{\sqrt{c}}\,
    K_{0}\!\Big(\tfrac{m}{\sqrt{c}}\sqrt{c\,D^{2}+b^{2}}\Big),
  \label{eq:K0}
\end{equation}
%
a Bessel $K_{0}$: the momentum-space Coulomb--Sturmian object, living on a
rational (genus-zero) curve.  This is why the two-center engine is weight one.

Two different scales cannot be linearized together.  Their product is
governed by the algebraic curve
%
\begin{equation}
  y^{2} = (c_{1}k^{2}+1)(c_{2}k^{2}+1),
  \label{eq:curve}
\end{equation}
%
a quartic with four distinct branch points $k=\pm i/\sqrt{c_{1}},\,
\pm i/\sqrt{c_{2}}$ whenever $c_{1}\neq c_{2}$: an \emph{elliptic curve},
genus one.  When $c_{1}=c_{2}$ the branch points collide and the quartic
becomes a perfect square---rational, genus zero.  The cleanest witness is the
period.  At $D_{1}=D_{2}=0$ the radial integral is exactly a complete
elliptic integral \textbf{[MEASURED, 31 digits; classical]},
%
\begin{equation}
  \int_{0}^{\infty}\!\frac{dk}{\sqrt{(c_{1}k^{2}+1)(c_{2}k^{2}+1)}}
  = \frac{1}{a_{1}\sqrt{c_{1}c_{2}}}\,K(m),
  \label{eq:period}
\end{equation}
%
with $a_{i}=1/\sqrt{c_{i}}$ and modulus $m=1-c_{\min}/c_{\max}$
(nondegenerate, $0<m<1$), while on the diagonal $c_{1}=c_{2}=c$ it collapses
to the elementary $\pi/(2\sqrt{c})$.  Since $c_{1}=s(1-s)$ and $c_{2}=t(1-t)$
range independently over $[0,\tfrac14]$, the full integral~\eqref{eq:reduced}
is an integral over a \emph{family} of elliptic curves, of modulus
$m(s,t)=1-c_{\min}/c_{\max}$, degenerate only on the measure-zero locus
$s=t$ or $s=1-t$.  Every term generated by the $\partial_{\zeta}$'s is a
meromorphic function on the same curve~\eqref{eq:curve}, so the whole integral
is a period of this elliptic family.

The structural statement is that the embedding tier of the transcendence
taxonomy (Papers~18,~34) is \emph{genus-graded}: genus zero
($\{E_{1},\ln,\gamma\}$) at two centers, genus one (elliptic) at three.
The two two-center negatives---no $\pi$, no dilogarithm---are genus-zero
properties; the third center does not climb to a higher rung of the
polylogarithm weight tower but leaves that tower for the elliptic one.  This
is the transcendence-level form of ``no coordinate system for three foci,''
and it also explains, after the fact, why an attempt to identify the three-center
value against dilogarithm constants was structurally doomed: the constants are
in the wrong class.

%% ========================================================================
\section{The reduction: an elliptic Bessel moment}
\label{sec:reduction}
%% ========================================================================

The radial factor can be named precisely.  Applying Eq.~\eqref{eq:K0} to each
factor and Parseval's theorem, the radial core (at fixed Feynman parameters,
before the elementary $j_{0}$ and $\partial_{\zeta}$ wrapping) is a two-mass
Bessel moment \textbf{[MEASURED, 30 digits]},
%
\begin{equation}
  M = \frac{2}{\pi\sqrt{c_{1}c_{2}}}
  \int_{0}^{\infty}\!\! db\;
  K_{0}\!\big(a_{1}\sqrt{p_{1}^{2}+b^{2}}\big)\,
  K_{0}\!\big(a_{2}\sqrt{p_{2}^{2}+b^{2}}\big),
  \label{eq:besselmoment}
\end{equation}
%
with $a_{i}=1/\sqrt{c_{i}}$ and $p_{i}=D_{i}\sqrt{c_{i}}$.  This is a
shifted-argument, two-scale member of the Bessel-moment family; it is not one
of the single-scale moments $\int_{0}^{\infty}t^{k}K_{0}^{n}(t)\,dt$
catalogued in closed form by Bailey, Borwein, Broadhurst and
Glasser~\cite{bbbg2008}.

A change of variable $x=\sqrt{c_{1}k^{2}+1}$ on the one-mass slice
$N(D)\equiv M(D,0)$ puts it in the cleanest canonical form: the Laplace
transform of the holomorphic elliptic differential,
%
\begin{equation}
  N(D) = \frac{1}{\sqrt{c_{1}}}
  \int_{1}^{\infty}\!
  \frac{e^{-Dx}\,dx}{\sqrt{(x^{2}-1)(\rho\,x^{2}+1-\rho)}},
  \qquad \rho=\frac{c_{2}}{c_{1}}.
  \label{eq:laplace}
\end{equation}
%
This is a single explicit one-dimensional integral of an elementary
integrand---already sharper than the infinite series of the standard
treatment~\cite{ozdogan2012}, and it makes the elliptic structure manifest.
It is an \emph{elliptic deformation of the Bessel} $K_{0}$: as $\rho\to0$
(coincident scales) $N\to K_{0}(D)$, and at $D=0$ it is the complete elliptic
integral $K$ of Eq.~\eqref{eq:period}.

Because $1/\sqrt{Q}$, $Q(x)=(x^{2}-1)(\rho x^{2}+1-\rho)$, satisfies the
first-order relation $Q f' + \tfrac12 Q' f = 0$, and the boundary at the
branch point $x=1$ vanishes ($Q(1)^{1/2}=0$), integration by parts against
$e^{-Dx}$ (using $\int_{1}^{\infty}x^{n}e^{-Dx}f\,dx=(-1)^{n}L^{(n)}(D)$,
$L\equiv\sqrt{c_{1}}\,N$) yields an exact fourth-order Picard--Fuchs equation
with polynomial coefficients \textbf{[SYMBOLIC PROOF; MEASURED residual
$\sim\!10^{-24}$]},
%
\begin{equation}
  D\rho\,L'''' + 2\rho\,L''' + D(1{-}2\rho)L'' + (1{-}2\rho)L'
  - D(1{-}\rho)L = 0.
  \label{eq:pf}
\end{equation}
%
Its characteristic polynomial factors as
$\rho(s^{2}-1)\big(s^{2}+(1-\rho)/\rho\big)$, whose roots $s=\pm1$ and
$s=\pm i\sqrt{(1-\rho)/\rho}$ are exactly the four branch points of the
curve~\eqref{eq:curve}: the $K_{0}/I_{0}$ sector from the $(x^{2}-1)$ pair,
the oscillatory $J_{0}/Y_{0}$ sector from the other.  A subordination
identity casts the same object as a two-scale sunrise integral,
$\propto\int\!\!\int ds_{1}ds_{2}\,(s_{1}s_{2})^{-1/2}
e^{-D_{1}^{2}/4s_{1}-s_{1}-D_{2}^{2}/4s_{2}-s_{2}}/\sqrt{s_{1}c_{1}+s_{2}c_{2}}$,
whose second-Symanzik-type $\sqrt{\text{linear form}}$ is the genus-one
signature and separates only when $c_{1}=c_{2}$.

One numerical fact pins the object's class beyond the period.  $N(D)$ is
\emph{non-analytic in $D$ at $D=0$}: it carries a $D\ln D$ term
\textbf{[MEASURED]}, so it is not a pure elliptic integral ($K$/$E$ are
analytic in their parameters).  The elliptic $D=0$ period together with the
$D\ln D$ nonanalyticity are the two fingerprints of an
\emph{elliptic ${\oplus}$ logarithmic} object---the inhomogeneous
(``regulator'') term of the Picard--Fuchs solution associated with the elliptic
dilogarithm~\cite{blochvanhove2015,adamsbognerweinzierl2014}---a
correspondence whose reach we test, and bound, in
Sec.~\ref{sec:obstruction}.

%% ========================================================================
\section{The elliptic-dilogarithm route is obstructed}
\label{sec:obstruction}
%% ========================================================================

The unequal-mass sunrise is closed by variation of parameters: a second-order
Picard--Fuchs operator in the external variable annihilates the two periods,
its inhomogeneity is an elementary subtopology (a tadpole), and the particular
solution---an iterated integral of the periods against that elementary
source---is the elliptic dilogarithm~\cite{adamsbognerweinzierl2014,
blochvanhove2015}.  Running the same machinery on $N(D)$ is the natural
attempt.  It does not close, and the reason is structural.

The curve $y^{2}=Q(x,\rho)$ of Eq.~\eqref{eq:curve} moves with the modulus
$\rho$, not with the Laplace variable $D$.  Accordingly the two periods
$K(\rho),K(1-\rho)$ are the homogeneous solutions of the second-order modulus
operator \textbf{[MEASURED: both annihilated to $10^{-18}$;
$L(0,\rho)=K(1-\rho)$ to $19$ digits]}
%
\begin{equation}
  \mathcal{M}_{\rho}
  = \rho(1-\rho)\,\partial_{\rho}^{2} + (1-2\rho)\,\partial_{\rho}
    - \tfrac14 ,
  \label{eq:modpf}
\end{equation}
%
and they are \emph{not} solutions of the fourth-order operator in
$D$: a constant gives $\mathcal{O}[1]=-D(1-\rho)\neq0$, and the single
Bessel functions of both sectors, $K_{0}(D),I_{0}(D)$ and
$J_{0}(\omega D),Y_{0}(\omega D)$ with $\omega=\sqrt{(1-\rho)/\rho}$, are
likewise not $D$-solutions \textbf{[MEASURED]} (they are coupled, which is why
Eq.~\eqref{eq:pf} does not factor into elementary Bessel operators).  The
elliptic-dilogarithm construction lives on the periods, so it must be run in
the modulus $\rho$.

There it meets the obstruction.  Acting with $\mathcal{M}_{\rho}$ under the
integral sign is an exact total derivative \textbf{[SYMBOLIC PROOF; exact
residual $0$]},
%
\begin{equation}
  \mathcal{M}_{\rho}\big[Q^{-1/2}\big]
  = \frac{d}{dx}\!\left[\,g(x)\,Q^{-1/2}\right],
  \quad
  g(x)=\frac{-x(x^{2}-1)}{4(\rho x^{2}+1-\rho)},
  \quad g(1)=0,
  \label{eq:exactform}
\end{equation}
%
so for $D\neq0$ the modulus source is
$S(D,\rho)\equiv\mathcal{M}_{\rho}[L]=D\!\int_{1}^{\infty}e^{-Dx}g\,Q^{-1/2}dx$.
Because $g$ \emph{vanishes at the branch point} $x=1$, differentiating in the
modulus produces no branch-point boundary term: the analogue of the sunrise's
pinched propagator---the elementary tadpole that would be the
inhomogeneity---is absent, and the source stays on the elliptic curve.
Quantitatively \textbf{[MEASURED, $45$ digits]}, $S$ closes exactly and
finitely only inside the integral's own Gauss--Manin $D$-module,
%
\begin{equation}
  S(D,\rho)=\sum_{k=0}^{3}P_{k}(D,\rho)\,\partial_{D}^{k}L,
  \qquad \deg_{D}P_{k}=2,
  \label{eq:inmodule}
\end{equation}
%
with held-out residual $5\times10^{-44}$, flat under increase of the fit
degree (numerically an exact finite identity); whereas the genuine
variation-of-parameters shape---period level $\{L,\partial_{D}L\}$ plus an
elementary Bessel subtopology $\{K_{0},K_{1}\}(D)$, $\{J_{0},Y_{0}\}(\omega
D)$---only \emph{approximates} it (residual $2\times10^{-14}$ falling
geometrically to $3\times10^{-30}$ across fit degrees, never saturating).  The
inhomogeneity is in-module, not a simpler subtopology, and variation of
parameters against the periods is therefore circular.

\textbf{[OBSERVATION]} $N(D)$ is thus one transcendence level above the
sunrise.  It is not an elliptic dilogarithm of elementary or period-level
arguments; it is a period of the rank-four irregular connection defined by
Eq.~\eqref{eq:pf}, for which the sunrise mechanism supplies no reduction to
the elliptic-polylogarithm tower.  Whether it admits a closed form of any
kind is a differential-Galois question---whether the fourth-order operator
factors over the differential field generated by the periods
$K(\rho),K(1-\rho)$---which we leave open.

%% ========================================================================
\section{Placement in the literature}
\label{sec:literature}
%% ========================================================================

The object of Eqs.~\eqref{eq:besselmoment}--\eqref{eq:pf} sits inside a
fully developed body of mathematics, applied so far only to problems far from
chemistry.  Two-mass elliptic Feynman integrals---the sunrise/banana
family---are evaluated in terms of elliptic dilogarithms as inhomogeneous
solutions of exactly this kind of Picard--Fuchs equation
\cite{blochvanhove2015,adamsbognerweinzierl2014,blochkerrvanhove2015}; their
periods are the complete elliptic integrals of Eq.~\eqref{eq:period}, and at
special (complex-multiplication) scale ratios their values reach critical
$L$-values of modular forms~\cite{broadhurst2016,fresansabbahyu2023}.  The
closest published analog to Eq.~\eqref{eq:besselmoment} is a doubled-argument
Bessel moment that Broadhurst evaluated as a product of complete elliptic
integrals via a lattice-Green-function contour~\cite{broadhurst2008}: the
same ``two-scale moment $\to$ product of $K$'s'' mechanism we expect at
rational $\rho$.

On the chemistry side the treatment is uniformly numerical: three-center
Slater repulsion integrals are reduced to infinite series or one-dimensional
quadrature~\cite{ozdogan2012,barnett_coulson1951}, and the Averys'
momentum-space Coulomb--Sturmian route~\cite{avery_avery_4center}, closest in
spirit to ours, likewise stops at expansion and numerics.  We could find no
paper drawing the bridge in either direction: no chemistry reference names an
electron-repulsion integral as a Bessel moment or an elliptic Feynman
integral, and no amplitudes reference applies the elliptic machinery to a
molecular integral.  \textbf{[OBSERVATION]} The identification---that the
three-center Slater integral, via the Fock projection, is a two-scale Bessel
moment on the sunrise elliptic curve---appears to be unmade.  This is an
absence from a targeted search, not a proof of absence; the residual risk is
an obscure single source.

%% ========================================================================
\section{Scope, and the open frontier}
\label{sec:scope}
%% ========================================================================

\emph{Scope.}  The demonstration is for $1s$ orbitals at a common exponent.
Higher $(n,\ell,m)$ multiply the integrand of Eq.~\eqref{eq:reduced} by
solid-harmonic polynomials in $\bm k$ and shift the exponents; these change
the weights and arguments but not the genus of the curve~\eqref{eq:curve},
which is set by the two distinct scales.  The result is a statement about the
transcendence \emph{class} of the three-center integral, established on its
simplest genuine instance.  It does not by itself deliver a polyatomic energy:
a molecular Hamiltonian requires the full three-center tensor, and a value
that is part closed-form and part numerical is neither more accurate (the
Gaussian fits are already exact to fit quality) nor decidable.  What the
reduction buys is an exact, fast, geometry-native evaluator for the
three-center integral and---the point of this paper---a precise diagnosis of
why it is hard.

\emph{The open problem.}  As shown in Sec.~\ref{sec:obstruction}, the
elliptic-dilogarithm route is obstructed: $N(D)$ is not an elliptic
dilogarithm of elementary or period-level data but a period of the rank-four
irregular connection~\eqref{eq:pf}.  The remaining question is whether that
connection admits \emph{any} closed form---equivalently, whether its
fourth-order operator factors over the differential field generated by the
periods $K(\rho),K(1-\rho)$---a differential-Galois problem we do not settle.
(This tempers a literature-scout expectation that the object would sit at the
tractable elliptic-dilogarithm level because it has only two propagator
factors; that pattern-match from the sunrise is overridden by the in-module
computation of Sec.~\ref{sec:obstruction}.)  Because it continues the Averys'
momentum-space Sturmian program and requires elliptic-Feynman-integral
technique one level beyond the sunrise, it is the natural place for
collaboration.

%% ========================================================================
\section{Conclusion}
%% ========================================================================

The three-center electron-repulsion integral is not merely a harder version
of the two-center one; it is an object of higher genus.  Two centers close at
genus zero over $\{E_{1},\ln,\gamma\}$; three centers, read in momentum space
where the coordinate obstruction dissolves, reduce to a two-scale Bessel
moment on an elliptic curve, an elliptic deformation of the Bessel $K_{0}$
with an explicit fourth-order Picard--Fuchs equation and a period that is a
complete elliptic integral.  The seventy-year difficulty of polyatomic Slater
integrals is, at the level of transcendence, a genus jump; and a classical
molecular integral turns out to be an instance of the modern theory of
elliptic Feynman integrals---a bridge that, to our knowledge, had not been
drawn.  The finite closed form remains open---and, we have shown, past the
reach of the sunrise's elliptic-dilogarithm mechanism, whose inhomogeneity is
absent here: the object is a period of a rank-four irregular connection, and
whether it closes at all is a differential-Galois question, posed here as the
frontier it deserves to be.

%% ========================================================================
\begin{thebibliography}{99}

\bibitem{fock1935}
V.~Fock, ``Zur Theorie des Wasserstoffatoms,''
\textit{Z.\ Phys.}\ \textbf{98}, 145--154 (1935).

\bibitem{avery1989}
J.~Avery, \textit{Hyperspherical Harmonics:\ Applications in Quantum
Theory} (Kluwer Academic, Dordrecht, 1989).

\bibitem{avery2006}
J.~E.~Avery and J.~S.~Avery, \textit{Generalized Sturmians and Atomic
Spectra} (World Scientific, Singapore, 2006).

\bibitem{avery_avery_4center}
J.~E.~Avery and J.~S.~Avery, ``4-Center STO interelectron repulsion
integrals with Coulomb Sturmians,'' \textit{Adv.\ Quantum Chem.}\ (Elsevier),
ScienceDirect S0065327617300710.

\bibitem{shibuya1965}
T.~Shibuya and C.~E.~Wulfman, ``Molecular orbitals in momentum space,''
\textit{Proc.\ R.\ Soc.\ Lond.\ A}\ \textbf{286}, 376 (1965).

\bibitem{ozdogan2012}
T.~\"Ozdo\u{g}an and E.~Ruiz, ``Evaluation of three-center two-electron
repulsion integrals over Slater-type orbitals,'' arXiv:1209.3755 (2012).

\bibitem{barnett_coulson1951}
M.~P.~Barnett and C.~A.~Coulson, ``The evaluation of integrals occurring
in the theory of molecular structure.\ Parts I \& II,''
\textit{Philos.\ Trans.\ R.\ Soc.\ Lond.\ A}\ \textbf{243}, 221--249 (1951).

\bibitem{fromm_hill1987}
D.~M.~Fromm and R.~N.~Hill, ``Analytic evaluation of the three-electron
integral,'' \textit{Phys.\ Rev.\ A}\ \textbf{36}, 1013--1044 (1987).

\bibitem{bbbg2008}
D.~H.~Bailey, J.~M.~Borwein, D.~Broadhurst, and M.~L.~Glasser, ``Elliptic
integral evaluations of Bessel moments and applications,''
\textit{J.\ Phys.\ A}\ \textbf{41}, 205203 (2008); arXiv:0801.0891.

\bibitem{broadhurst2008}
D.~Broadhurst, ``Elliptic integral evaluation of a Bessel moment by contour
integration of a lattice Green function,'' arXiv:0801.4813 (2008).

\bibitem{blochvanhove2015}
S.~Bloch and P.~Vanhove, ``The elliptic dilogarithm for the sunset graph,''
\textit{J.\ Number Theory}\ \textbf{148}, 328--364 (2015); arXiv:1309.5865.

\bibitem{adamsbognerweinzierl2014}
L.~Adams, C.~Bogner, and S.~Weinzierl, ``The two-loop sunrise graph in two
space-time dimensions with arbitrary masses in terms of elliptic
dilogarithms,'' \textit{J.\ Math.\ Phys.}\ \textbf{55}, 102301 (2014);
arXiv:1405.5640.

\bibitem{blochkerrvanhove2015}
S.~Bloch, M.~Kerr, and P.~Vanhove, ``A Feynman integral via higher normal
functions,'' \textit{Compos.\ Math.}\ \textbf{151}, 2329--2375 (2015);
arXiv:1406.2664.

\bibitem{broadhurst2016}
D.~Broadhurst, ``Feynman integrals, $L$-series and Kloosterman moments,''
arXiv:1604.03057 (2016).

\bibitem{fresansabbahyu2023}
J.~Fres\'an, C.~Sabbah, and J.-D.~Yu, ``Quadratic relations between Bessel
moments,'' \textit{Algebra Number Theory}\ \textbf{17}, 541--602 (2023);
arXiv:2006.02702.

\bibitem{loutey_paper18}
J.~Loutey, ``Exchange Constants: A Transcendence Taxonomy for the Geometric
Vacuum,'' Paper~18 in this series.

\bibitem{loutey_paper34}
J.~Loutey, ``A Projection Taxonomy for Physical Observables,'' Paper~34 in
this series.

\bibitem{loutey_paper58}
J.~Loutey, ``Angular Sparsity Is an Atomic-Sector Property: The Abelian
Residue of Multi-Center Bonding,'' Paper~58 in this series.

\bibitem{test_routeC}
\textsc{GeoVac} regression suite, \texttt{tests/test\_routeC\_momentum.py}
(momentum reproduction of the three-center integral; the single-dispersion
$K_{0}$ identity; the elliptic-curve period), and
\texttt{tests/test\_two\_center\_eri\_aabb.py} (two-center weight-one census).

\end{thebibliography}

\end{document}
